% Source PDF: source/普通高考/2011/2011辽宁理.pdf, page 1, question 6.
% Vector program flowchart; original PNG is retained at img/2011/liaoning_science/q6_fig1.png.
% Node -> directed-edge list: 开始 -> 输入 n -> s=0,t=1,k=1,p=1 -> k<n;
% 是 -> p=s+t -> s=t,t=p -> k=k+1 -> return to the test-entry stem;
% 否 -> output p -> 结束. The no branch runs right, down, then left into
% the output stem exactly as in the source; its single arrowhead is at the
% output-stem merge. All node text is locally zero-indented and centered.

\begin{tikzpicture}[
  >=Stealth,
  x=1cm, y=1cm,
  line width=0.45pt,
  line cap=round,
  line join=round,
  font=\normalsize,
  flowtext/.style={anchor=center,align=center,inner sep=0pt,
    execute at begin node={\setlength{\parindent}{0pt}}},
  flowbox/.style={flowtext,draw,minimum height=.68cm,inner xsep=3pt,inner ysep=2pt},
  terminal/.style={flowbox,rounded corners=2pt,minimum width=1.35cm},
  process/.style={flowbox},
  io/.style={flowbox,shape=trapezium,trapezium left angle=72,trapezium right angle=108,
    minimum height=.76cm},
  decision/.style={flowbox,shape=diamond,aspect=2.8,minimum width=3.70cm,
    minimum height=1.02cm,inner sep=0pt},
  branchlabel/.style={flowtext,inner sep=0pt}
]
  \node[terminal] (start) at (0,10.80) {开始};
  \node[io,minimum width=1.80cm] (input) at (0,9.35) {输入 $n$};
  \node[process,minimum width=5.25cm] (init) at (0,7.85)
    {$s=0,t=1,k=1,p=1$};
  \node[decision] (test) at (0,6.25) {$k<n$};
  \node[process,minimum width=2.30cm] (sum) at (0,4.80) {$p=s+t$};
  \node[process,minimum width=2.90cm] (shift) at (0,3.35) {$s=t,t=p$};
  \node[process,minimum width=2.35cm] (inc) at (0,1.90) {$k=k+1$};
  \node[io,minimum width=2.25cm] (output) at (0,0.25) {输出 $p$};
  \node[terminal] (finish) at (0,-1.15) {结束};

  \draw[->] (start.south) -- (input.north);
  \draw[->] (input.south) -- (init.north);
  \draw[->] (init.south) -- (test.north);

  \draw[->] (test.south) -- node[branchlabel,right=2pt] {是} (sum.north);
  \draw[->] (sum.south) -- (shift.north);
  \draw[->] (shift.south) -- (inc.north);

  % 是 branch loops to the connector immediately above the decision.
  \coordinate (loop-bottom) at (0,0.85);
  \coordinate (loop-left) at (-2.60,0.85);
  \coordinate (loop-entry) at (-2.60,6.70);
  \draw[->] (inc.south) -- (loop-bottom) -- (loop-left) -- (loop-entry) -- (0,6.70);

  % 否 branch runs independently right and down, then joins the output stem.
  \coordinate (merge) at (0,0.85);
  \draw (test.east) -- node[branchlabel,above=2pt] {否} (3.45,6.25);
  \draw (3.45,6.25) -- (3.45,0.85) -- (merge);
  \draw[->] (merge) -- (output.north);
  \draw[->] (output.south) -- (finish.north);

  \path[use as bounding box] (-2.95,-1.50) rectangle (3.85,11.10);
\end{tikzpicture}
